
\subsection{同伦等价与算子 $\mathfrak q$ I：改变内部闭上链}
\label{subsec:4-6-4}
\liftIndex{I178}

在第 4.6 节余下的小节中，我们用第 3.8 节引入的算子
$\mathfrak q$，讨论闭上链
\[
\mathfrak b\in C^2(M;\Lambda^+_{0,\mathrm{nov}})
=C(M;\Lambda^+_{0,\mathrm{nov}})[2]^0
\]
对滤过 $A_\infty$ 代数
$(C(L;\Lambda_{0,\mathrm{nov}}),\mathfrak m)$
的形变。第 3.8.5 小节已定义形变代数
$(C(L;\Lambda_{0,\mathrm{nov}}),\mathfrak m^{(\mathfrak b,0)})$；
以下简记为 $(C(L;\Lambda_{0,\mathrm{nov}}),
\mathfrak m^{\mathfrak b})$。
本小节证明：在严格同伦等价的意义下，这个代数
只依赖于 $\mathfrak b$ 的上同调类
（定理 4.6.47）。

作为对照，对任意余链
$b\in C^1(L;\Lambda^+_{0,\mathrm{nov}})$，
我们在第 3.6.1 小节定义了
$(C(L;\Lambda_{0,\mathrm{nov}}),\mathfrak m^b)$
（见命题 3.6.10）。
它在通常的、非严格的同伦等价意义下不依赖于 $b$；
也就是说，甚至不依赖于 $b$ 的规范等价类
（注 4.3.17）。

本小节的论证纯属代数性质。
结论对任何具有满足定理 3.8.32 结论的算子
$\mathfrak q$ 的滤过 $A_\infty$ 代数 $C$ 均成立，
但我们只按几何例子陈述。
以下三个小节（4.6.4--4.6.6）假设
$R\supseteq\mathbb Q$，并沿用第 3.8 节的记号。

\begin{theorem}{4.6.47}\label{thm:4-6-47}
设 $\mathfrak b_0,\mathfrak b_1
\in C^2(M;\Lambda^+_{0,\mathrm{nov}})$
是代表 $H^2(M;\Lambda_{0,\mathrm{nov}})$ 中
同一个上同调类的闭上链。
则滤过 $A_\infty$ 代数
$(C,\mathfrak m^{\mathfrak b_0})$
与 $(C,\mathfrak m^{\mathfrak b_1})$
严格同伦等价。
\end{theorem}

定理 4.6.47 在第 7.4.6 小节证明；
具体见定理 7.4.132\translatornote{源书把证明定位到定理 7.4.118；本书在定理 7.4.132 后明确写明定理 4.6.47 由定理 Y 与该定理推出，故改正回指。见疑误 A-0591。}。
我们还需要下面的命题，
它将在定义 4.6.50 中使用。
设 $\mathfrak b_i
\in C^2(M;\Lambda^+_{0,\mathrm{nov}})$
（$i=0,1,2$）为代表同一个上同调类的闭上链
\footnote{译者注：原书此处未写次数 $2$，
且称它们代表同一个“同调类”；
但所用定理 4.6.47 要求二次闭上链及其上同调类。
见疑误 A-0247。}。令
\begin{equation}
\mathfrak f_{ij}\colon
(C,\mathfrak m^{\mathfrak b_i})
\longrightarrow
(C,\mathfrak m^{\mathfrak b_j})
\tag{4.6.48}\label{eq:4-6-48}
\end{equation}
为定理 4.6.47 给出的同伦等价。
\begin{proposition}{4.6.49}\label{prop:4-6-49}
$\mathfrak f_{12}\circ\mathfrak f_{01}$
与 $\mathfrak f_{02}$ 同伦\translatornote{源书写作未定义的 ``shu-homotopic''；本命题在命题 7.4.134 中重述并证明时只写 homotopic，故不补造修饰语。见疑误 A-0593。}。
\end{proposition}
命题 4.6.49 也在第 7.4.6 小节证明；
具体见命题 7.4.134\translatornote{源书误称 ``Theorem 7.4.134''；该编号在后文明确属于命题。见疑误 A-0592。}。

由命题 4.6.49，可将定义 4.3.26 推广如下。
\begin{statement}{定义--推论}{4.6.50}\label{defcor:4-6-50}
设 $(\mathfrak b_0,b_0),(\mathfrak b_1,b_1)
\in\widehat{\mathcal M}_{\mathrm{def}}(L)$。
若
\[
[\mathfrak b_0]=[\mathfrak b_1]
\in H^2(M;\Lambda_{0,\mathrm{nov}})
\]
且 $\mathfrak f_*(b_0)\sim b_1$，
就称这两对数据彼此规范等价，
并记为
$(\mathfrak b_0,b_0)\sim(\mathfrak b_1,b_1)$
\liftIndex{I179}。
\footnote{译者注：原书把
$[\mathfrak b_0]=[\mathfrak b_1]$
置于 $H^1(L;\Lambda_{0,\mathrm{nov}})$；
但 $\mathfrak b_i$ 是 $M$ 上的二次闭上链，
此处所需正是定理 4.6.47 的环境上同调类条件。
见疑误 A-0248。}
这里
\[
\mathfrak f\colon
(C(L;\Lambda_{0,\mathrm{nov}}),
 \mathfrak m^{\mathfrak b_0})
\longrightarrow
(C(L;\Lambda_{0,\mathrm{nov}}),
 \mathfrak m^{\mathfrak b_1})
\]
是定理 4.6.47 中的同伦等价。

命题 4.6.49 保证这种规范等价是等价关系。
同样可在
$\widehat{\mathcal M}_{\mathrm{weak,def}}(L)$
上定义等价关系 $\sim$。
$\widehat{\mathcal M}_{\mathrm{def}}(L)$、
$\widehat{\mathcal M}_{\mathrm{weak,def}}(L)$
关于 $\sim$ 的商集，
分别与定义 4.3.26 中的
$\mathcal M_{\mathrm{def}}(L)$、
$\mathcal M_{\mathrm{weak,def}}(L)$ 相同。
\end{statement}

\begin{lemma}{4.6.51}\label{lem:4-6-51}
若 $(\mathfrak b_i,b_i)
\in\widehat{\mathcal M}_{\mathrm{weak,def}}(L)$
且 $(\mathfrak b_0,b_0)\sim(\mathfrak b_1,b_1)$，
则
\[
\mathfrak{PO}(\mathfrak b_0,b_0)
=\mathfrak{PO}(\mathfrak b_1,b_1).
\]
\end{lemma}
证明见第 7.4.6 小节。

有自然投影
$\pi_{\mathrm{amb}}\colon
\mathcal M_{\mathrm{def}}(L)
\to H^2(M;\Lambda^+_{0,\mathrm{nov}})$。
由引理 4.6.51，势函数
\[
\mathfrak{PO}_i\colon
\mathcal M_{\mathrm{weak,def}}(L^{(i)})
\longrightarrow\Lambda_{0,\mathrm{nov}}^{(0)}
\]
与投影
$\pi_{\mathrm{amb},i}\colon
\mathcal M_{\mathrm{weak,def}}(L^{(i)})
\to H^2(M;\Lambda^+_{0,\mathrm{nov}})$
均有定义；参见第 4.3.3 小节。
这里
\[
\begin{aligned}
\mathcal M_{\mathrm{weak,def}}(L^{(i)})
=\bigl\{(\mathfrak b,b)\ \bigm|\;&
\mathfrak b\in H^2(M;\Lambda^+_{0,\mathrm{nov}}),\\
&b\in\mathcal M_{\mathrm{weak}}
\bigl(C(L^{(i)};\Lambda_{0,\mathrm{nov}}),
      \mathfrak m^{\mathfrak b}\bigr)\bigr\}.
\end{aligned}
\]
此处对 $\mathfrak b$ 的上同调类，
使用第 4.3.3 小节选定的代表元理解形变结构。

第 5.2.4 小节的命题 5.2.37 (3)、5.2.38 (2)
将证明：在给定规范等价类内改变定界余链时，
第 3.8.7 小节定义的弗洛尔上同调不变。
因此
\[
HF\bigl((L^{(1)},(\mathfrak b_1,b_1)),
        (L^{(0)},(\mathfrak b_0,b_0));
        \Lambda_{0,\mathrm{nov}}\bigr)
\]
由下列数据参数化：
\[
\begin{aligned}
&\bigl([\mathfrak b_1,b_1],
       [\mathfrak b_0,b_0]\bigr)\\
&\quad\in
\mathcal M_{\mathrm{weak,def}}(L^{(1)})
\times_{\,
 H^2(M;\Lambda^+_{0,\mathrm{nov}})
 \times\Lambda_{0,\mathrm{nov}}^{(0)}}
\mathcal M_{\mathrm{weak,def}}(L^{(0)})\\
&\quad=
\bigl\{
([\mathfrak b_1,b_1],[\mathfrak b_0,b_0])
\ \bigm|\
\pi_{\mathrm{amb},1}[\mathfrak b_1,b_1]
=\pi_{\mathrm{amb},0}[\mathfrak b_0,b_0],\\
&\hspace{33mm}
\mathfrak{PO}_1[\mathfrak b_1,b_1]
=\mathfrak{PO}_0[\mathfrak b_0,b_0]
\bigr\}.
\end{aligned}
\]
下文也用同一记号表示
$(\mathfrak b_i,b_i)$ 依定义 4.6.50 所得的规范等价类。
